\chapter*{Transformatoren und Schaltgruppen}
\label{chap:04-transformatoren}

Dieses Kapitel ist als Bestandteil des modularen Lehrskripts vorgesehen.
Inhalte werden ergänzt, nachdem die fachlichen Grundlagen und die
Quellenlage für dieses Thema systematisch ausgearbeitet wurden.

\begin{center}
\begin{minipage}{0.8\textwidth}
\noindent
\begin{tikzpicture}[transform shape]
	% Paths, nodes and wires:
	\draw (2, 4) to[european inductor, l={$L_{s1}=\sigma_1 L_1$}] (4, 4);
	\draw (0, 4) to[generic, l={$R_1$}] (2, 4);
	\draw (5, 2) to[european inductor, l_={$L_{h1}$}, label distance=0.03cm] (5, 4);
	\draw (10, 4) to[generic, l={$R_2$}] (12, 4);
	\draw (8, 4) to[european inductor, l={$L_{s2}=\sigma_2 L_2$}] (10, 4);
	\draw (4, 2) to[generic, l={$R_{FE}$}] (4, 4);
	\node[transformer core, cute, xscale=0.95, yscale=0.95] at (7.003, 3.003){};
	\draw (6.005, 4) |- (5, 4);
	\draw (6, 2) |- (5, 2) |- (4, 2) -- (0, 2);
	\draw (4, 4) -- (5, 4);
	\draw (8, 2.005) -| (12, 2);
	\draw[line width=1pt, -latex] (0, 3.75) -- (0, 2.25)   node[midway, right=2pt] {$\underline{U}_1$};
	\draw[line width=1pt, -latex] (12, 3.75) -- (12, 2.25) node[midway, right=2pt] {$\underline{U}_2$};
	\draw[line width=1pt, -latex] (6, 3.75) -- (6, 2.25)   ;%node[midway, right=2pt] {$\underline{U}_{1h}$};
	\draw[line width=1pt, -latex] (8, 3.75) -- (8, 2.25)   node[midway, right=2pt] {$\underline{U}_{2h} = \frac{N_2}{N_1} \, \underline{U}_{1h} $};
    \node[above=3pt] at (5.8, 4) {$L_{h1} = (1-\sigma_1) L_1 $};
\end{tikzpicture}
\begin{tikzpicture}[transform shape]
	% Paths, nodes and wires:
	\draw (2, 4) to[european inductor, l={$L_{s1}=\sigma_1 L_1$}] (4, 4);
	\draw (0, 4) to[generic, l={$R_1$}] (2, 4);
	\draw (5, 2) to[european inductor, l_={$L_{h1}$}, label distance=0.03cm] (5, 4);
	\draw (10, 4) to[generic, l={$R_2$}] (12, 4);
	\draw (8, 4) to[european inductor, l={$L_{s2}=\sigma_2 L_2$}] (10, 4);
	\draw (4, 2) to[generic, l={$R_{FE}$}] (4, 4);
	\node[transformer core, cute, xscale=0.95, yscale=0.95] at (7.003, 3.003){};
	\draw (6.005, 4) |- (5, 4);
	\draw (6, 2) |- (5, 2) |- (4, 2) -- (0, 2);
	\draw (4, 4) -- (5, 4);
	\draw (8, 2.005) -| (12, 2);
	\draw[line width=1pt, -latex] (0, 3.75) -- (0, 2.25)   node[midway, right=2pt] {$\underline{U}_1$};
	\draw[line width=1pt, -latex] (12, 3.75) -- (12, 2.25) node[midway, right=2pt] {$\underline{U}_2$};
	\draw[line width=1pt, -latex] (6, 3.75) -- (6, 2.25)   ;%node[midway, right=2pt] {$\underline{U}_{1h}$};
	\draw[line width=1pt, -latex] (8, 3.75) -- (8, 2.25)   node[midway, right=2pt] {$\underline{U}_{2h} = \frac{N_2}{N_1} \, \underline{U}_{1h} $};
    \node[above=3pt] at (5.8, 4) {$L_{h1} = (1-\sigma_1) L_1 $};
\end{tikzpicture}
\par\bigskip
\noindent
\begin{tikzpicture}[transform shape]
	% Paths, nodes and wires:
	\draw (2, 4) to[european inductor, l={$L_{s1}=\sigma_1 L_1$}] (4, 4);
	\draw (0, 4) to[generic, l={$R_1$}] (2, 4);
	\draw (5, 2) to[european inductor, l_={$L_{h1}$}, label distance=0.03cm] (5, 4);
	\draw (8.005, 4) to[generic, l={$R_2'$}, i^<=$i_2'$] (10.005, 4);
	\draw (6.005, 4) to[european inductor, l={$L_{s2}'=\sigma_2 L_2'$}] (8.005, 4);
	\draw (4, 2) to[generic, l={$R_{FE}$}] (4, 4);
	\draw (6.005, 4) |- (5, 4);
	\draw (10, 2) -- (5, 2) |- (4, 2) -- (0, 2);
	\draw (4, 4) -- (5, 4);
    \draw (6, 2) |- (5, 2) |- (4, 2) -- (0, 2);
	\draw (8, 2.005) -| (10, 2);
	\draw (2, 4) to[european inductor, l={$L_{s1}=\sigma_1 L_1$}] (4, 4);
	\draw (0, 4) to[generic, l={$R_1$}] (2, 4);
	\draw (5, 2) to[european inductor, l_={$L_{h1}$}, label distance=0.03cm] (5, 4);
	\draw (4, 2) to[generic, l={$R_{FE}$}] (4, 4);
	\node[transformer core, cute, xscale=0.95, yscale=0.95] at (10.997, 3.003){};
	\draw[line width=1pt, -latex] (0, 3.75) -- (0, 2.25)   node[midway, right=2pt] {$\underline{U}_1$};
	\draw[line width=1pt, -latex] (12, 3.75) -- (12, 2.25) node[midway, right=2pt] {$\underline{U}_2$};
	\draw[line width=1pt, -latex] (10, 3.75) -- (10, 2.25) node[midway, left=0pt] {$\frac{N_1}{N_2} \, \underline{U}_2$};
	\draw[line width=1pt, -latex] (6, 3.75) -- (6, 2.25)  node[midway, right=2pt]  {$\underline{U}_{1h} $};
\end{tikzpicture}
\par\bigskip
\noindent
\begin{tikzpicture}[transform shape]
	% Paths, nodes and wires:
	\draw (2, 4) to[european inductor, l={$L_{s1}=\sigma_1 L_1$}] (4, 4);
	\draw (0, 4) to[generic, l={$R_1$}] (2, 4);
	\draw (5, 2) to[european inductor, l_={$L_{h1}$}, label distance=0.03cm] (5, 4);
	\draw (8, 4) to[generic, l={$R_2'$}, i^<=$i_2'$]  (10, 4);
	\draw (6, 4) to[european inductor, l={$L_{s2}'$}] (8, 4);
	\draw (4, 2) to[generic, l={$R_{FE}$}] (4, 4);
	\draw (6, 4) |- (5, 4);
	\draw (10, 2) -- (5, 2) |- (4, 2) -- (0, 2);
	\draw (4, 4) -- (5, 4);
	\draw[line width=1pt, -latex] (0, 3.75) -- (0, 2.25) node[midway, right=2pt] {$\underline{U}_1$};
	\draw[line width=1pt, -latex] (10, 3.75) -- (10, 2.25) node[midway, right=2pt] {$\underline{U}_2'$};
	\draw[line width=1pt, -latex] (6, 3.75) -- (6, 2.25) node[midway, right=2pt] {$\underline{U}_{1h}$};
\end{tikzpicture}

\par\bigskip
\noindent
\begin{tikzpicture}[transform shape]
	% Paths, nodes and wires:
	\draw (4.5, 5)    to[european inductor, l={$X_{1\sigma}+X_{2\sigma}'$}] (6.5, 5);
	\draw (2.5, 5)    to[generic, l={$R_1+R_2'$}] (4.5, 5);
	\draw (1.5, 2.5)  to[european inductor, l={$2X_{1h}$}, label distance=0.03cm] (1.5, 4.5);
	\draw (2.5, 2.5)  to[generic, l_={$2R_{FE}$}] (2.5, 4.5);
	\draw (6.5, 2.5)  to[european inductor, l={$2X_{1h}$}, label distance=0.03cm] (6.5, 4.5);
	\draw (7.5, 2.5)  to[generic, l_={$2R_{FE}$}] (7.5, 4.5);
    \draw[line width=1pt, -latex] (0, 4.5) -- (0, 2.5) node[midway, left=2pt] {$\underline{U}_{1}$};
	\draw[line width=1pt, -latex] (9, 4.5) -- (9, 2.5) node[midway, right=2pt] {$\underline{U}_{2}'$};
	\draw (1.5, 2.5) -- (2.5, 2.5);
	\draw (1.5, 4.5) -- (2.5, 4.5);
	\draw (0, 5) -- (2.5, 5);
    \draw (0, 5) to[short, i^>={$I_1$}] (1,5);
	\draw (2, 4.5) -| (2, 5);
	\draw (6.5, 2.5) -- (7.5, 2.5);
	\draw (6.5, 4.5) -- (7.5, 4.5);
	\draw (7, 4.5) -| (7, 5);
	\draw (0, 2) -- (9, 2);
	\draw (2, 2.5) -| (2, 2);
	\draw (7, 2.5) -| (7, 2);
	\draw (6.5, 5) -- (9, 5);
    \draw (8, 5) to[short, i^<={$I_2'$}] (9,5);
\end{tikzpicture}


\begin{tikzpicture}[transform shape]
	\draw (1, 0.75) to[tmultiwire] (1, 0.25);
	\draw (1, 2.5) to[ioosource] (1, -0);
	\draw (1, 2.25) to[tmultiwire] (1, 1.75);
\end{tikzpicture}

\end{minipage}
\end{center}
$$ L_1 = L_{h1} + L_{s1} $$
$$ L_{h1} = (1-\sigma_1) L_1 $$

$$ M = \frac{N_1}{N_2} = k \, \sqrt{L_1 L_2} $$


\chapter{Transformatoren und ihre Ersatzschaltbilder}
\label{chap:transformatoren}

\section{Lernziele und Einordnung}

Transformatoren ermöglichen die verlustarme Übertragung elektrischer Energie
zwischen unterschiedlichen Spannungsebenen. Für die Netzberechnung sind sie
nicht nur als ideale Spannungsübersetzer zu betrachten. Wicklungswiderstände,
Streuflüsse, Magnetisierungsstrom, Eisenverluste und die Schaltgruppe bestimmen
das Betriebsverhalten und beeinflussen insbesondere Kurzschluss- und Erdschlussströme.

Nach diesem Kapitel sollen folgende Aufgaben gelöst werden können:
\begin{itemize}
  \item Übersetzungsverhältnis und Nennströme aus Typenschilddaten bestimmen;
  \item ein einphasiges Ersatzschaltbild eines Drehstromtransformators aufstellen;
  \item Leerlauf- und Kurzschlussversuch auswerten;
  \item aus der Kurzschlussspannung eine Ersatzimpedanz bestimmen;
  \item Wirk- und Blindanteil der Transformatorimpedanz unterscheiden;
  \item die Auswirkungen von Schaltgruppe und Neutralpunkterdung auf das
        Nullsystem qualitativ beurteilen.
\end{itemize}

\section{Idealer Transformator}

\subsection{Induktionsgesetz und Spannungsübersetzung}

Die elektromagnetische Induktion wird durch das Faradaysche
Induktionsgesetz beschrieben. In differentieller Form der
Maxwellgleichungen enthalten als
\[
    \nabla \times \vec{E}
    = -\frac{\partial \vec{B}}{\partial t}.
\]
Sie beschreibt, wie eine zeitlich veränderliche magnetische
Flussdichte $\vec{B}$ ein elektrisches Wirbelfeld erzeugt.
Durch Integration über eine von
einer Wicklung umschlossene Fläche folgt das Induktionsgesetz für
eine einzelne Windung. Für eine Wicklung mit $N$ Windungen ergibt
sich unter Berücksichtigung der gewählten Bezugsrichtungen
\[
    u(t) = N \frac{\mathrm d\Phi(t)}{\mathrm dt},
\]
wobei $\Phi$ der magnetische Fluss durch eine Windungsfläche ist.
Das Vorzeichen hängt von der Festlegung der Wicklungsrichtung bezüglich
der positiven Flussrichtung ab.

\subsection{Entstehung des magnetischen Flusses}

Der magnetische Fluss im Transformatorkern wird durch den Strom
in den Wicklungen hervorgerufen. Das Durchflutungsgesetz, oder
in vollständiger, differentieller Form, das
Ampère-Maxwell-Gesetz

\[
        \nabla \times \vec{H}
        = \vec{J} + \frac{\partial}{\partial t} \vec{D}
\]

beschreibt mit $\vec{H}$ die entstehende magnetische Feldstärke
aufgrund von elektrischer Stromdichte $\vec{J}$ und in seiner
Wirkung als Verschiebungsstrom bezeichneten Ableitung der elektrischen Flussdichte $\vec{D}$. Das
Gesetz besagt, dass sowohl elektrische Ströme als auch zeitlich veränderliche
elektrische Felder ein magnetisches Wirbelfeld hervorrufen.

Der Wicklungsstrom $N I$ bildet für einen Trafo die Durchflutung,
die nach diesem Gesetz mit einer mittleren
magnetischen Weglänge $l_\mathrm{m}$ dem entstehenden magnetischen Feld gleicht.

\[
    H l_\mathrm{m} = N I.
\]

Die magnetische Flussdichte ergibt sich aus der Materialgleichung
\[
    \vec{B} = \mu_0 \mu_\mathrm{r}\vec{H},
\]
wobei $\mu_0$ die magnetische Feldkonstante und $\mu_\mathrm{r}$
die relative Permeabilität des Kernmaterials bezeichnet. Der
magnetische Fluss durch den Kernquerschnitt $A_\mathrm{K}$ ist
\[
    \Phi = \int_{A_\mathrm{K}} \vec{B}\cdot \mathrm d\vec{A}.
\]
Bei näherungsweise homogener Flussdichte und einem konstanten
Kernquerschnitt vereinfacht sich dies zu
\[
    \Phi \approx B A_\mathrm{K}.
\]

Für den nachfolgend skizzierten Transformator ist die als $\underline{U}_1$
eingespeiste Spannung die Ursache für einen Strom $i_1$, der zu einem Fluss
$\Phi$ führt und eine der Ursache entgegenwirkende, induszierte Spannung
aufbaut, die das Treiben von Strom zu dem sich einstellenden Arbeitspunkt begrenzt.

\vspace{-1.1em}
\begin{center}
\begin{tikzpicture}[transform shape]
	% Paths, nodes and wires:
	\node[shape=rectangle, fill={rgb,255:red,120;green,150;blue,165}, draw, line width=1pt, minimum width=4.965cm, minimum height=5.715cm] at (3.49, 2.86){};
	\draw[line width=1pt, dash pattern={on 4pt off 2pt on 1pt off 2pt}, -latex] (3.5, 5) -- (5, 5) -- (5, 0.75) -- (2, 0.75) -- (2, 5) -- (3.75, 5) node[midway, above=2pt] {$\underline{\Phi}(t)$};
	\node[shape=rectangle, fill={rgb,255:red,255;green,255;blue,255}, draw, line width=1pt, minimum width=2.477cm, minimum height=3.715cm] at (3.496, 2.86){};
	\node[shape=rectangle, fill={rgb,255:red,209;green,209;blue,209}, draw, line width=1pt, minimum width=2.251cm, minimum height=3.215cm] at (1.597, 2.86){};
	\node[shape=circle, draw, line width=1pt, minimum width=0.465cm] at (2.404, 3.604){};
	\draw[line width=1pt] (2.227, 3.78) -- (2.581, 3.427);
	\draw[line width=1pt] (2.227, 3.427) -- (2.581, 3.78);
	\node[shape=circle, draw, line width=1pt, minimum width=0.465cm] at (0.802, 2.485){};
	\node[circ]  at (0.802, 2.485){};
	\node[shape=circle, draw, line width=1pt, minimum width=0.465cm] at (0.802, 3.985){};
	\node[circ]  at (0.802, 3.985){};
	\node[shape=circle, draw, line width=1pt, minimum width=0.465cm] at (0.802, 1.735){};
	\node[circ]  at (0.802, 1.735){};
	\node[shape=circle, draw, line width=1pt, minimum width=0.465cm] at (2.404, 2.854){};
	\draw[line width=1pt] (2.227, 3.03) -- (2.581, 2.677);
	\draw[line width=1pt] (2.227, 2.677) -- (2.581, 3.03);
	\node[shape=circle, draw, line width=1pt, minimum width=0.465cm] at (2.404, 2.104){};
	\draw[line width=1pt] (2.227, 2.28) -- (2.581, 1.927);
	\draw[line width=1pt] (2.227, 1.927) -- (2.581, 2.28);
	\node[shape=circle, draw, line width=1pt, minimum width=0.465cm] at (0.802, 3.235){};
	\node[circ]  at (0.802, 3.235){};
	\node[shape=rectangle, xscale=-1, yscale=-1, fill={rgb,255:red,209;green,209;blue,209}, draw, line width=1pt, minimum width=2.251cm, minimum height=3.215cm] at (5.37, 2.875){};
	\node[shape=circle, xscale=-1, yscale=-1, draw, line width=1pt, minimum width=0.465cm] at (6.177, 2.808){};
	\node[circ]  at (6.177, 2.808){};
	\node[shape=circle, xscale=-1, yscale=-1, draw, line width=1pt, minimum width=0.465cm] at (6.177, 2.058){};
	\node[circ]  at (6.177, 2.058){};
	\node[shape=circle, xscale=-1, yscale=-1, draw, line width=1pt, minimum width=0.465cm] at (4.575, 3.177){};
	\draw[line width=1pt] (4.399, 3.354) -- (4.752, 3);
	\draw[line width=1pt] (4.399, 3) -- (4.752, 3.354);
	\node[shape=circle, xscale=-1, yscale=-1, draw, line width=1pt, minimum width=0.465cm] at (4.575, 2.427){};
	\draw[line width=1pt] (4.399, 2.604) -- (4.752, 2.25);
	\draw[line width=1pt] (4.399, 2.25) -- (4.752, 2.604);
	\node[shape=circle, xscale=-1, yscale=-1, draw, line width=1pt, minimum width=0.465cm] at (6.177, 3.558){};
	\node[circ]  at (6.177, 3.558){};
	\draw (-1.25, 4) to[sinusoidal voltage source] (-1.25, 1.75);
	\draw (-1.25, 4) -| (0.802, 3.985);
	\draw (-1.25, 1.75) -| (0.802, 1.735); % node[anchor=south] at ([yshift=-0.12cm]N1.north){$i_1$};
	\draw (7.5, 3.564) -| (6.177, 3.56); %\node[flowarrow](N1) at (-0.25, 4.25){} node[anchor=south] at ([yshift=-0.12cm]N1.north){$i_1$};
	\draw (7.5, 2.048) -| (6.221, 2.062);
	\draw[line width=1pt, -latex] (7.25, 3.266) -- (7.25, 2.266) node[midway, right=2pt] {$\underline{U}_2$};
	\draw[line width=1pt, -latex] (0, 3.5) -- (0, 2.25) node[midway, left=2pt] {$\underline{U}_1$};
	\node[flowarrow](N1) at (-0.5, 4){} node[anchor=south] at ([yshift=-0.12cm]N1.north){$i_1$};
	\node[flowarrow, xscale=-1](N2) at (7.154, 3.567){} node[anchor=south] at ([yshift=-0.12cm]N2.north){$i_2$};
\end{tikzpicture}
\end{center}

Damit folgt unmittelbar für den Leerlauf ($i_2 = 0\,A$) eine
dem Windungsverhältnis proportional induzierte Spannung $\underline{U}_2$
sowie für den Kurzschlussfall ein dem Windungsverhältnis antiproportionaler
Strom.

Die Gleichungen beschreiben den Zusammenhang zwischen
Wicklungsstrom, magnetischer Feldstärke, Flussdichte und
magnetischem Fluss. Im realen Transformator wird der Fluss durch
den nichtlinearen Kernwerkstoff, den magnetischen Widerstand,
Streuanteile und gegebenenfalls die Sättigung beeinflusst. Für
den idealen Transformator wird angenommen, dass der gesamte
magnetische Fluss beide Wicklungen vollständig verkettet.

\subsection{Herleitung des Übersetzungsverhältnisses}

Primär- und Sekundärwicklung sind auf demselben magnetischen
Kreis angeordnet. Im idealen Modell verkettet deshalb derselbe
zeitabhängige Fluss $\Phi(t)$ beide Wicklungen. Wendet man das
Induktionsgesetz auf beide Wicklungen an, erhält man betragsmäßig
\[
    U_1 = N_1 \left|\frac{\mathrm d\Phi}{\mathrm dt}\right|,
    \qquad
    U_2 = N_2 \left|\frac{\mathrm d\Phi}{\mathrm dt}\right|.
\]
Das Verhältnis der beiden Spannungen ergibt sich unmittelbar,
weil der gemeinsame Flussterm herausgekürzt werden kann:
\[
    \boxed{\frac{U_1}{U_2}=\frac{N_1}{N_2}=a.}
\]
Das Übersetzungsverhältnis $a$ ist hier als Verhältnis der
Primärspannung zur Sekundärspannung beziehungsweise der
Primärwindungszahl zur Sekundärwindungszahl definiert. Bei dieser
Konvention gilt $a>1$ für einen abwärts und $a<1$ für einen
aufwärts übersetzenden Transformator.

Bei sinusförmigem magnetischem Fluss
\[
    \Phi(t)=\Phi_{\max}\sin(\omega t)
\]
ist die induzierte Spannung ebenfalls sinusförmig. Aus der
Ableitung des Flusses folgt für den Effektivwert
\[
    \boxed{U=4{,}44\,f\,N\,\Phi_{\max}.}
\]
Die Beziehung gilt für sinusförmigen Fluss und vernachlässigte
Spannungsabfälle an den Wicklungswiderständen und Streuinduktivitäten.

\subsection{Stromübersetzung und Leistungsübertragung}

Für den idealen Transformator werden Wicklungswiderstände,
Streuinduktivitäten und Kernverluste vernachlässigt. Es wird
daher keine Wirkleistung in Wärme umgesetzt. Im sinusförmigen
stationären Betrieb gilt für die Beträge der Wirkleistung bei
idealer Leistungsübertragung
\[
    U_1 I_1 = U_2 I_2.
\]
Zusammen mit dem Spannungsübersetzungsverhältnis folgt
\[
    \boxed{\frac{I_1}{I_2}=\frac{U_2}{U_1}=\frac{1}{a}.}
\]
Die Stromübersetzung verhält sich somit umgekehrt zur
Spannungsübersetzung: Eine höhere Sekundärspannung geht mit
einem entsprechend geringeren Sekundärstrom einher, sofern die
übertragene Leistung gleich bleibt.

Die Strombeziehung ist eine Betragsbeziehung. Die Vorzeichen
beziehungsweise Phasenlagen der Ströme hängen von den gewählten
Bezugsrichtungen und der Punktkennzeichnung der Wicklungen ab.


Die induzierte Spannung einer Wicklung folgt aus dem Induktionsgesetz von
Faraday, in differenzieller Form der Maxwellgleichung:

\[
    %\boxed{
        \nabla \times \vec{E}
        = -\frac{\partial \vec{B}}{\partial t}
    %}
\]

durch Integration über eine Windung und Summation über $N$ Windungen zu

\[
  u(t)=N\frac{\mathrm d\Phi(t)}{\mathrm dt},
\]
wobei das Vorzeichen von der gewählten Wicklungs- und Flussrichtung abhängt.

Bei sinusförmigem magnetischem Fluss
\[
  \Phi(t)=\Phi_{\max}\sin(\omega t)
\]
ergibt sich für den Effektivwert der induzierten Spannung
\[
  U=4{,}44\,f\,N\,\Phi_{\max}.
\]

Für zwei ideal gekoppelte Wicklungen folgt daraus
\[
  \frac{U_1}{U_2}=\frac{N_1}{N_2}=a.
\]
Das Übersetzungsverhältnis \(a\) ist hier als Verhältnis der Spannungen
auf der Seite 1 zur Seite 2 definiert.

Bei idealer Leistungsübertragung gilt
\[
  U_1 I_1=U_2 I_2
\]
für die Beträge der entsprechenden Effektivwerte. Damit folgt
\[
  \frac{I_1}{I_2}=\frac{1}{a}.
\]

Bei einem realen Drehstromtransformator müssen die Leiter-Leiter-Spannungen,
Strangspannungen, Wicklungsverschaltungen und gegebenenfalls die Phasenverschiebung
zwischen den Seiten getrennt betrachtet werden. Das einfache Verhältnis \(a\)
ist unmittelbar auf die Spannungen der jeweils gekoppelten Wicklungen anwendbar.

\subsection{Nennströme}

Für einen symmetrischen Drehstromtransformator mit Bemessungsscheinleistung
\(S_N\) und Leiter-Leiter-Bemessungsspannung \(U_N\) beträgt der Bemessungsstrom
auf der betreffenden Seite
\[
  I_N=\frac{S_N}{\sqrt{3}\,U_N}.
\]

Beispiel: Ein Transformator mit
\[
  S_N=40\,\mathrm{MVA},\qquad U_{N,1}=110\,\mathrm{kV}
\]
besitzt auf der 110-kV-Seite den Bemessungsstrom
\[
  I_{N,1}
  =\frac{40\,\mathrm{MVA}}{\sqrt{3}\cdot110\,\mathrm{kV}}
  \approx210\,\mathrm{A}.
\]

Bei \(U_{N,2}=20\,\mathrm{kV}\) ergibt sich sekundärseitig
\[
  I_{N,2}
  =\frac{40\,\mathrm{MVA}}{\sqrt{3}\cdot20\,\mathrm{kV}}
  \approx1{,}155\,\mathrm{kA}.
\]

Die Bemessungsleistung bleibt näherungsweise gleich, während sich die
Bemessungsströme umgekehrt zum Spannungsverhältnis verhalten.

\section{Reales Ersatzschaltbild}

\subsection{Serienzweig}

Das reale Ersatzschaltbild berücksichtigt die Wicklungswiderstände
\(R_1\) und \(R_2\) sowie die Streureaktanzen \(X_{\sigma 1}\) und
\(X_{\sigma 2}\).

Die Wicklungswiderstände beschreiben die ohmschen Verluste der Leiter.
Die Streureaktanzen repräsentieren den Anteil des magnetischen Flusses,
der nicht beide Wicklungen in gleicher Weise verkettet.

Für die Sekundärimpedanz gilt beim Bezug auf die Primärseite
\[
  Z_2'=a^2 Z_2.
\]
Daraus folgen für den zusammengefassten Serienzweig näherungsweise
\[
  R_k=R_1+a^2R_2,
  \qquad
  X_k=X_{\sigma 1}+a^2X_{\sigma 2},
\]
und
\[
  Z_k=R_k+jX_k.
\]

Diese Beziehungen setzen voraus, dass alle Größen auf dieselbe Seite
bezogen sind. Die Streuinduktivitäten sind dabei durch ihre Reaktanzen
bei der betrachteten Frequenz repräsentiert:
\[
  X_\sigma=\omega L_\sigma.
\]

\subsection{Magnetisierungszweig}

Der Magnetisierungszweig wird durch einen Eisenverlustwiderstand
\(R_{\mathrm{Fe}}\) und eine Magnetisierungsreaktanz \(X_m\) modelliert.
Der Widerstand repräsentiert die im Kern auftretenden Verluste,
insbesondere Hysterese- und Wirbelstromverluste. Die Reaktanz bildet
den zur Erzeugung des magnetischen Flusses erforderlichen Strom ab.

Das vereinfachte Ersatzschaltbild enthält somit:
\begin{itemize}
  \item den idealen Transformator;
  \item den Serienwiderstand und die Streureaktanz;
  \item den parallel liegenden Magnetisierungszweig.
\end{itemize}

Für Kurzschlussrechnungen wird der Magnetisierungszweig häufig vernachlässigt,
weil der Kurzschlussstrom hauptsächlich durch die Serienimpedanz begrenzt wird.
Für Leerlauf, Einschaltvorgänge und genauere Betriebsberechnungen ist diese
Vereinfachung nicht ohne Weiteres zulässig.

\section{Leerlaufversuch}

\subsection{Messprinzip}

Beim Leerlaufversuch wird eine Wicklung mit einer definierten Spannung gespeist,
während die andere Wicklung offen bleibt. Gemessen werden beispielsweise
\[
  U_0,\qquad I_0,\qquad P_0.
\]

Bei Nennfrequenz und geeigneter Prüfspannung ist der Leerlaufstrom typischerweise
deutlich kleiner als der Bemessungsstrom. Deshalb sind die ohmschen Verluste
der gespeisten Wicklung häufig klein gegenüber den Eisenverlusten.

Unter dieser Näherung gilt
\[
  P_{\mathrm{Fe}}\approx P_0.
\]

Die Näherung wird schlechter, wenn die Wicklungsverluste des Leerlaufstroms
nicht mehr vernachlässigbar sind oder die Prüfbedingungen von den
Bemessungsbedingungen abweichen.

\subsection{Bestimmung der Ersatzparameter}

Für einen einphasigen Ersatzschaltkreis ergibt sich näherungsweise
\[
  R_{\mathrm{Fe}}=\frac{U_0^2}{P_0}.
\]

Der Wirkstromanteil des Leerlaufstroms ist
\[
  I_{\mathrm{Fe}}=\frac{P_0}{U_0}.
\]

Der Magnetisierungsstrom wird bei angenommener Orthogonalität von Wirk- und
Blindstromanteil bestimmt durch
\[
  I_m=\sqrt{I_0^2-I_{\mathrm{Fe}}^2}.
\]
Damit folgt
\[
  X_m=\frac{U_0}{I_m}.
\]

Bei Drehstromtransformatoren ist zu beachten, ob die Messwerte als
Leitergrößen oder als Stranggrößen vorliegen. Die Parameterformeln sind
auf den jeweils verwendeten einphasigen Ersatzkreis anzuwenden.

\section{Kurzschlussversuch}

\subsection{Messprinzip}

Beim Kurzschlussversuch wird eine Wicklung kurzgeschlossen. An die andere
Wicklung wird eine reduzierte Spannung angelegt, sodass der gewünschte
Prüfstrom fließt. Für den üblichen Versuch mit Bemessungsstrom werden
Spannung, Strom und Wirkleistung gemessen:
\[
  U_k,\qquad I_k,\qquad P_k.
\]

Da die angelegte Spannung deutlich unter der Bemessungsspannung liegt,
ist der magnetische Fluss geringer als im Nennbetrieb. Die Eisenverluste
sind deshalb meist klein gegenüber den Wicklungsverlusten.

Die gemessene Kurzschlussleistung kann daher näherungsweise zur Bestimmung
der Serienverluste verwendet werden.

\subsection{Auswertung für einen einphasigen Ersatzkreis}

Für den einphasigen Ersatzkreis gelten
\[
  |Z_k|=\frac{U_k}{I_k},
\]
\[
  R_k=\frac{P_k}{I_k^2},
\]
und
\[
  X_k=\sqrt{|Z_k|^2-R_k^2}.
\]

Die letzte Gleichung setzt voraus, dass \(P_k\) die Wirkleistung des
betrachteten einphasigen Ersatzkreises beschreibt. Bei Drehstrommessungen
muss die Gesamtleistung konsistent auf den Strang bezogen werden.

Die Kurzschlussimpedanz lautet anschließend
\[
  Z_k=R_k+jX_k.
\]

Eine rein betragsmäßige Angabe der Kurzschlussimpedanz reicht für viele
Netzberechnungen als Näherung aus. Für genauere Berechnungen, insbesondere
bei niedrigen Spannungen oder hohem Widerstandsanteil, ist die Trennung
zwischen \(R_k\) und \(X_k\) wichtig.

\section{Kurzschlussspannung und Impedanz aus Typenschilddaten}

Die relative Kurzschlussspannung wird häufig in Prozent angegeben:
\[
  u_k=\frac{U_k}{U_N}\cdot100\%.
\]

Dabei beziehen sich \(U_k\) und \(U_N\) auf dieselbe Wicklungsseite und
dieselbe Spannungsdefinition. Bei symmetrischen Drehstromtransformatoren
kann aus der relativen Kurzschlussspannung der Betrag der Strangimpedanz
näherungsweise bestimmt werden:
\[
  |Z_k|=\frac{u_k}{100}\frac{U_N^2}{S_N}.
\]

In dieser Gleichung ist \(U_N\) die Leiter-Leiter-Bemessungsspannung
und \(S_N\) die dreiphasige Bemessungsscheinleistung. Die Gleichung liefert
die Impedanz des einphasigen Strangersatzes.

\subsection{Beispiel}

Gegeben sei ein Transformator mit
\[
  S_N=10\,\mathrm{MVA},\quad
  U_N=20\,\mathrm{kV},\quad
  u_k=8\%.
\]

Damit ergibt sich
\[
  |Z_k|
  =0{,}08\frac{(20\,\mathrm{kV})^2}{10\,\mathrm{MVA}}
  =3{,}2\,\Omega.
\]

Der Bemessungsstrom auf der 20-kV-Seite beträgt
\[
  I_N=\frac{10\,\mathrm{MVA}}{\sqrt{3}\cdot20\,\mathrm{kV}}
  \approx289\,\mathrm{A}.
\]

Der Betrag der Kurzschlussimpedanz ist damit bestimmt. Die Aufteilung
in Wirk- und Blindanteil ist aus \(u_k\) allein jedoch nicht möglich.
Dafür sind weitere Herstellerangaben oder Messdaten erforderlich.

\section{Stufenschalter und Spannungsregelung}

Stufenschalter verändern das Übersetzungsverhältnis des Transformators.
Bei einer Wicklungsanzapfung wird die wirksame Windungszahl und damit
die Übersetzung verändert. Je nach Bauart erfolgt die Umschaltung
spannungsfrei oder unter Last.

Ein Stufenschalter dient dazu, die Spannung auf der geregelten Seite
innerhalb des vorgesehenen Bereichs zu halten. Die Wirkung hängt von
Netzimpedanzen, Laststrom, Leistungsfaktor und der gewählten Regelstrategie ab.

Eine Spannungsänderung ist nicht mit einer Veränderung der Kurzschlussimpedanz
gleichzusetzen. Bei detaillierten Kurzschlussrechnungen müssen jedoch die
tatsächliche Stufenstellung und die für das Berechnungsverfahren geltenden
Vorgaben berücksichtigt werden.

\section{Schaltgruppen und Nullsystem}

Die Schaltgruppe beschreibt die Wicklungsverschaltung und die relative
Phasenverschiebung zwischen den Seiten. Sie ist für Lastfluss, Parallelbetrieb
und Fehlerberechnung von Bedeutung.

Bei der Analyse müssen mindestens folgende Eigenschaften unterschieden werden:
\begin{itemize}
  \item Phasenverschiebung zwischen Ober- und Unterspannungsseite;
  \item zugänglicher Neutralpunkt und dessen Erdung;
  \item Möglichkeit von Nullsystemströmen in den Wicklungen;
  \item magnetischer Rückweg des Nullsystemflusses.
\end{itemize}

Eine Dreieckwicklung kann einen Umlaufweg für bestimmte Nullsystemströme
bereitstellen. Daraus folgt aber nicht automatisch, dass Nullsystemströme
zwischen den beiden Netzseiten übertragen werden. Das Verhalten ist anhand
der konkreten Schaltgruppe, der Wicklungsanordnung und der Erdung zu beurteilen.

\section{Parallelbetrieb}

Für den Parallelbetrieb von Transformatoren sind insbesondere zu prüfen:
\begin{itemize}
  \item gleiches oder kompatibles Übersetzungsverhältnis;
  \item passende Phasenverschiebung und Schaltgruppe;
  \item gleiche Phasenfolge und richtige Klemmenzuordnung;
  \item vergleichbare Kurzschlussspannungen;
  \item geeignete Impedanzwinkel und Bemessungsleistungen.
\end{itemize}

Abweichende Übersetzungsverhältnisse können Ausgleichsströme verursachen.
Abweichende Kurzschlussimpedanzen führen zu einer ungleichen Aufteilung
der Lastströme. Ein Parallelbetrieb darf daher nicht allein aufgrund
gleicher Nennspannungen angenommen werden.

\section{Zusammenfassung}

Für Netzberechnungen ist die Transformatorimpedanz ein zentrales Element.
Die Kurzschlussspannung liefert deren Betrag auf der gewählten Bezugsseite.
Die Trennung in Wirk- und Blindanteil erfordert zusätzliche Informationen.
Schaltgruppe und Erdung sind besonders für Nullsystem- und Erdschlussmodelle
entscheidend.

\section{Übungsaufgaben}

\begin{enumerate}
  \item Berechnen Sie die Bemessungsströme eines Transformators mit
        \(S_N=25\,\mathrm{MVA}\) und \(U_N=110/20\,\mathrm{kV}\).
  \item Ein Transformator besitzt \(S_N=16\,\mathrm{MVA}\),
        \(U_N=10\,\mathrm{kV}\) und \(u_k=6\%\).
        Bestimmen Sie den Betrag der Strang-Kurzschlussimpedanz auf der
        10-kV-Seite.
  \item Erläutern Sie, warum sich aus \(u_k\) allein nicht bestimmen lässt,
        welcher Anteil der Kurzschlussimpedanz ohmsch ist.
  \item Beschreiben Sie, welche Informationen zur Bewertung des
        Nullsystemverhaltens eines Transformators erforderlich sind.
\end{enumerate}
